\section*{Proof of the small core lemma: lemma \ref{lem:small_core}}
Recall the statement of Lemma \ref{lem:small_core}:

\medskip

    For every $\vocab$-structure $\cG$, for every $\ell \in \bbN$,
    $\cG$ has a $\ell$-core $\cH$ with size at most $(2^{2n+4m+2}) \cdot \ell$.

\medskip


\begin{proof}
    Recall from the body that the intuition behind the construction is to repeatedly select a 1-type or tuple that violates the $\ell$-core condition
    and add the relevant vertices to $\cH$.
    Since there are only finitely many such choices, the process eventually stabilizes,
    and the size of the constructed core remains bounded.

    We formalize this intuition with the following $2^{2n+4m+1}$ round procedure.
    Let $\cH_0$ be an empty $\vocab$-structure.
    For the $i^{th}$ round,
    we select either:
    \begin{compactitem}
        \item
            a 1-type $\pi$ that is realized in $\cG \setminus \cH_{i-1}$ by at most $\ell$ vertices, or

        \item
            a tuple $\tuple{\pi_1, \eta, \pi_2} \in \onetypes \times \twotypes \times \onetypes$ that is realized in $\cG \setminus \cH_{i-1}$ by at most $\ell$ pairs of vertices.
    \end{compactitem}
    We define $\cH_{i}$ as follows:
    \begin{compactitem}
        \item
           If a 1-type $\pi$ is selected, then $\cH_i$ is obtained by extending $\cH_{i-1}$
           with all vertices from $\cG \setminus \cH$ with type $\pi$. 
            Note that there are at most $\ell$ vertices added.
        \item
            If a tuple $\tuple{\pi_1, \eta, \pi_2}$ is selected,
            then $\cH_i$ is obtained by extending  $\cH_{i-1}$ with all vertices from $\cG \setminus \cH$ involved in a pair realizing $\tuple{\pi_1, \eta, \pi_2}$.
            Note that there are at most $2\ell$ vertices added.
        \item Otherwise if there is no such 1-type of tuple, $\cH_i := \cH_{i-1}$.
    \end{compactitem}
    Finally, $\cH := \cH_{2^{2n+4m+1}}$.

    Note that in the $i^{th}$ round, we add at most $2\ell$ vertices to $\cH_i$.
    Thus the size of $\cH$ is bounded by $2^{2n+4m+1} \cdot 2\ell = 2^{2n+4m+2} \cdot \ell$.
    We claim that $\cH$ is an $\ell$-core of $\cG$.
    In the $i^{th}$ round, if a 1-type $\pi$ is selected,
    then for every $j \ge i$,
    no vertex realized $\pi$ remains in $\cG \setminus \cH_j$.
    A similar property holds if a tuple is selected.
    Since there are $2^{n+m}$ possible 1-types and $2^{2n+4m}$ possible tuples,
    after $2^{2n+4m+1} \ge 2^{n+m} + 2^{2n+4m}$ rounds,
    no valid selection remains.
    Hence, $\cH$ is an $\ell$-core of $\cG$.
\end{proof}